\section{Lo que enseña o1: RL, CoT y el pattern de reflexión}

La sección anterior señaló las deficiencias de la next token prediction; esta sección examina cómo entiende o1 Zhang Xiangyu (张祥雨). En la entrevista, su juicio es claro: lo esencial de o1 no es solo que use RL, sino que activa un pattern de cadena de pensamiento. Los nombres de los numerosos algoritmos del RL tradicional no son lo importante; lo importante es que el modelo ya vio durante el preentrenamiento muchos patrones de pensar, reflexionar, verificar y rehacer, y el RL puede reforzarlos.

\begin{figure}[H]
\centering
\includegraphics[width=0.96\textwidth]{rl-cot-pattern.png}
\caption{RL y CoT Pattern: el RL no es solo el nombre de un algoritmo; lo esencial es activar patrones de pensamiento en los que se pueda buscar. Diagrama conceptual propio, elaborado a partir de la conversación de 00:50:48--01:01:52.}
\end{figure}

\begin{knowledgebox}{Lectura de la figura: el RL no es magia; el pattern es la palanca}
En la figura, Pretrain aporta los patrones existentes, el RL los refuerza con retroalimentación sobre el objetivo, CoT despliega los pasos, Reflection corrige la trayectoria y Critical Decision resuelve las bifurcaciones decisivas. Lo que enseña o1 es que el modelo no aprende a pensar de la nada, sino que amplifica los patterns de pensamiento que ya existían, dispersos, en el preentrenamiento.
\end{knowledgebox}

\subsection{Critical Decision: donde un solo token no basta}

Esta subsección explica por qué importa la reflexión. Durante el razonamiento, no todos los tokens son igual de decisivos. Muchos tokens solo dan formato a la expresión; lo que realmente determina el éxito o el fracaso son unas pocas critical decisions: ir a la izquierda o a la derecha, qué fórmula usar, si hace falta volver atrás a verificar, si una conclusión intermedia es confiable. Cuando la complejidad de elegir una rama supera la capacidad expresiva de un solo token, el modelo necesita desplegar más pasos y descomponer el problema de decisión en un espacio en el que se pueda buscar.

\begin{figure}[H]
\centering
\includegraphics[width=0.96\textwidth]{meta-cot-critical-decision.png}
\caption{Meta-CoT y Critical Decision: cuando un solo token no basta para elegir una rama, se necesitan la reflexión y una metacadena de pensamiento. Diagrama conceptual propio, elaborado a partir de la conversación de 01:01:52--01:10:57.}
\end{figure}

\begin{knowledgebox}{Lectura de la figura: Meta-CoT es el CoT del CoT}
De State a Branch, y después a One token?, Reflect y Meta-CoT, la figura expresa «volver a pensar sobre el proceso de pensamiento». El modelo no solo escribe pasos, sino que, en los nodos decisivos, juzga si la trayectoria actual basta, si debe retroceder o si debe cambiar de rama.
\end{knowledgebox}

\begin{importantbox}{Nota de clase: el espacio de acciones importa más que el nombre del algoritmo}
Zhang Xiangyu insiste en que muchas veces lo determinante no es lo misterioso que sea un algoritmo de RL concreto, sino si el espacio de acciones contiene la acción correcta. Sin acciones de reflexión, al modelo le cuesta buscar en los puntos de decisión críticos; con acciones de reflexión, el RL tiene algo que reforzar.
\end{importantbox}

\subsection{Resumen de la sección}

Lo que o1 enseña a la multimodalidad es que la inteligencia no consiste solo en modelos más grandes: también requiere un proceso intermedio en el que se pueda buscar. CoT, la reflexión y Meta-CoT convierten decisiones de un solo token, que no se pueden resolver, en un espacio de acciones que puede descomponerse, explorarse por ensayo y error, y optimizarse.
